\section{Proofs of the main results}\label{app:proofs}

\subsection{Proof of Lemma~\ref{lem:exact_correction}}
\begin{proof}
Fix $\theta$ and $z\in\mathbb R^d$, and put $C=U+iV$.
The fundamental theorem of calculus and the exponential envelope give, for every $q\in\mathbb C^d$,
\begin{equation}\label{eq:appendix_score_envelope}
\|g^{\mathbb C}(\theta;q)-g(\theta;0)\|
\le M(\theta)\|q\|e^{a\|q\|}.
\end{equation}
Since $C$ is complex Gaussian, its norm has exponential moments of every linear order. Thus $g^{\mathbb C}(\theta;z+C)$ is absolutely integrable. The same bound permits exchange of the expectation and angular integral below.

For every $\phi\in[0,2\pi]$, the real pair
\[
(\operatorname{Re}(e^{i\phi}C),\operatorname{Im}(e^{i\phi}C))
=(U\cos\phi-V\sin\phi,\;U\sin\phi+V\cos\phi)
\]
has the same centered Gaussian law as $(U,V)$: its covariance is $\operatorname{diag}(\Omega,\Omega)$. Hence $e^{i\phi}C\overset{d}=C$. On the other hand, for each fixed $c\in\mathbb C^d$, the one-variable function $s\mapsto g^{\mathbb C}(\theta;z+sc)$ is entire. Cauchy's mean-value formula, applied componentwise, gives
\[
\frac{1}{2\pi}\int_0^{2\pi}
g^{\mathbb C}(\theta;z+e^{i\phi}c)\,d\phi
=g^{\mathbb C}(\theta;z)=g(\theta;z).
\]
Taking expectation over $C$ and using rotational invariance yields
\[
\mathbb E\,g^{\mathbb C}(\theta;z+C)
=\mathbb E\!\left[\frac{1}{2\pi}\int_0^{2\pi}
g^{\mathbb C}(\theta;z+e^{i\phi}C)\,d\phi\right]
=g(\theta;z).
\]
Taking real parts proves the first claim. For the second, condition on $\mathcal F_{t-1}$, apply the first claim conditional on $Z_t$, and then average over $Z_t$. Predictability and the assumed interchange of gradient and expectation give $\nabla F(\theta_t)$.
\end{proof}

\subsection{Reduction to classical SGD noise conditions}
\begin{proof}
Put $D=U+iV$. The fundamental theorem of calculus gives
\[
\|\operatorname{Re}g^{\mathbb C}(\theta;Z+D)-g(\theta;Z)\|
\le M(\theta)\|D\|e^{a(\|Z\|+\|D\|)}.
\]
Independence of $Z$ and $D$ gives
$\mathbb E\|\Delta(\theta)\|^2\le K_aM(\theta)^2$.
Equation~\eqref{eq:amplitude_control} then gives the correction-error
bound in the main text.
Sub-Gaussian vectors in finite dimension have exponential moments of their
norm of every linear order; the same holds for the Gaussian pair $(U,V)$,
regardless of whether $\Omega$ is singular. Hence $K_a<\infty$.
Conditioning first on $Z_t$ and using exact correction shows that
$\Delta_t$ is orthogonal in $L^2$ to
$g(\theta_t;Z_t)-\nabla F(\theta_t)$. This proves
\eqref{eq:relative_variance}. In particular,
\begin{equation}\label{eq:score_second_moment}
\mathbb E_t\|G_t\|^2
\le\sigma^2+(1+\rho)\|\nabla F(\theta_t)\|^2.
\end{equation}
For the CLTs, stronger moments follow from the same input bound.
Set $b_a(q)=\|q\|e^{a\|q\|}$ and $Q=Z+U+iV$. Centering cancels
$g(\theta;0)$ and gives
\[
\|G(\theta)-\nabla F(\theta)\|
\le M(\theta)\bigl[b_a(Q)+\mathbb E b_a(Z)\bigr].
\]
The bracket has moments of every finite order. Hence, for every $r\ge2$,
\begin{equation}\label{eq:higher_noise_moments}
\mathbb E_t\|\xi_t\|^r\le C_r M(\theta_t)^r
\le C_r\bigl(A+B\|\nabla F(\theta_t)\|^2\bigr)^{r/2},
\end{equation}
where $C_r$ is deterministic and finite. Under the strong convexity and
smoothness conditions used for the CLTs, almost-sure convergence gives
\[
\sup_t\|\nabla F(\theta_t)\|<\infty\quad\text{almost surely}.
\]
Thus the conditional higher moments are bounded along each convergent
trajectory by a finite random constant. On a fixed neighborhood of
$\theta^*$ the bound is deterministic, since
$\|\nabla F(\theta)\|\le L\|\theta-\theta^*\|$.
These bounds supply the Lindeberg control below without a global
uniform bound on $M$ or a moment assumption on the corrected score.
\end{proof}

\subsection{Convex convergence}
\begin{proof}
Write $e_t=\theta_t-\theta^*$ and $f_t=F(\theta_t)-F(\theta^*)$.
For a convex, $L$-smooth function,
\[
\langle\nabla F(\theta_t),e_t\rangle
\ge f_t+\frac1{2L}\|\nabla F(\theta_t)\|^2.
\]
Expanding the squared distance, using unbiasedness and
\eqref{eq:score_second_moment}, gives
\[
\mathbb E_t\|e_{t+1}\|^2
\le\|e_t\|^2-2\eta_tf_t
+\bigl((1+\rho)\eta_t^2-\eta_t/L\bigr)
\|\nabla F(\theta_t)\|^2+\sigma^2\eta_t^2.
\]
The gradient term is nonpositive under the stated step-size bound.
Taking expectations, summing, and applying convexity to the weighted
average proves the displayed bound. The rates follow from
$\sum_t\eta_t^2<\infty$ and the growth of $\sum_{t\le T}\eta_t$.
\end{proof}

\subsection{Strong convexity and MSE}
\begin{proof}
Write $e_t=\theta_t-\theta^*$. Strong convexity and smoothness give
$\langle\nabla F(\theta_t),e_t\rangle\ge\mu\|e_t\|^2$ and
$\|\nabla F(\theta_t)\|\le L\|e_t\|$. Hence
\begin{equation}\label{eq:strong_recursion}
\mathbb E_t\|e_{t+1}\|^2
\le\bigl(1-2\mu\eta_t+(1+\rho)L^2\eta_t^2\bigr)\|e_t\|^2
+\sigma^2\eta_t^2.
\end{equation}
This inequality first ensures finite second moments at every finite time.
The Robbins--Siegmund theorem \cite{robbinssiegmund1971}, applied with
summable coefficients $(1+\rho)L^2\eta_t^2$ and $\sigma^2\eta_t^2$,
then shows that $\|e_t\|^2$ converges almost surely and
$\sum_t\eta_t\|e_t\|^2<\infty$. Since $\sum_t\eta_t=\infty$,
the limit is zero.

For the rates, put $u_t=\mathbb E\|e_t\|^2$, $h_t=t+\gamma$,
$A=2\mu c$, $B=(1+\rho)L^2c^2$, and $D=\sigma^2c^2$.
Taking expectations gives
\begin{equation}\label{eq:mse_scalar}
u_{t+1}\le(1-Ah_t^{-\alpha}+Bh_t^{-2\alpha})u_t
+Dh_t^{-2\alpha}.
\end{equation}
If $\alpha<1$, choose $N$ so that the coefficient is nonnegative and
at most $1-(A/2)h_t^{-\alpha}$ for $t\ge N$.
Choose $M\ge4D/A$ with $u_N\le Mh_N^{-\alpha}$.
Increasing $N$ if needed, $h_t^{1-\alpha}\ge4\alpha/A$.
The induction hypothesis $u_t\le Mh_t^{-\alpha}$ then implies
\[
u_{t+1}\le Mh_t^{-\alpha}-\frac{AM}{4}h_t^{-2\alpha}
\le M(h_t+1)^{-\alpha},
\]
where the last inequality uses
$h^{-\alpha}-(h+1)^{-\alpha}\le\alpha h^{-\alpha-1}$.
After adjusting $M$ at this fixed starting index, induction proves
$u_t=O(t^{-\alpha})$.

If $\alpha=1$, choose $N$ so that $q_j=1-A/h_j+B/h_j^2>0$ for
$j\ge N$. Since $\log q_j=-A/h_j+O(h_j^{-2})$,
\[
\prod_{j=s}^{t-1}q_j\le C(h_s/h_t)^A\qquad(t>s\ge N),
\]
with a constant independent of $s,t$. Iterating \eqref{eq:mse_scalar}
therefore gives
\[
u_t\le C h_t^{-A}
\left(h_N^Au_N+D\sum_{s=N}^{t-1}h_s^{A-2}\right).
\]
The sum is $O(t^{A-1})$ for $A>1$, $O(\log t)$ for $A=1$,
and $O(1)$ for $A<1$. These are the three asserted rates, and all
converge to zero.
\end{proof}

\subsection{Preparatory estimates for the central limit theorems}
\label{app:clt_preparation}
Throughout the CLT proofs, write
$\bar e_T=T^{-1}\sum_{t=1}^Te_t=\bar\theta_T-\theta^*$,
and put $h=\lambda_{\min}(H)>0$ and
\[
r(e)=\nabla F(\theta^*+e)-He.
\]
Local Lipschitz continuity of the Hessian gives a radius $\varepsilon_0>0$
and a constant $C_r$ such that
$\|r(e)\|\le C_r\|e\|^2$ for $\|e\|\le\varepsilon_0$.
Globally, smoothness gives $\|r(e)\|\le(L+\|H\|)\|e\|$.

\paragraph{Localization.}
Choose $0<\varepsilon\le\varepsilon_0$ and define
$\tau_N=\inf\{t\ge N:\|e_t\|>\varepsilon\}$, with
$\inf\varnothing=\infty$. Almost-sure convergence implies
\begin{equation}\label{eq:localization_probability}
\mathbb P(\tau_N<\infty)\longrightarrow0\qquad(N\to\infty).
\end{equation}
For $\alpha<1$, the global MSE bound yields
\[
\mathbb E[\|r(e_t)\|\mathbf1_{\{\tau_N>t\}}]
\le C\,t^{-\alpha}.
\]
For $\alpha=1$, the local spectral condition $ch>1/2$ permits a choice
$m\in(1/(2c),h)$ and a sufficiently small $\varepsilon$ such that
$\langle\nabla F(\theta^*+e),e\rangle\ge m\|e\|^2$
on this ball. The event $\{\tau_N>t\}$ is measurable with respect to
$\mathcal F_{t-1}$. Setting
$v_t=\mathbb E[\|e_t\|^2\mathbf1_{\{\tau_N>t\}}]$, discarding
paths that exit at the next step gives
\[
v_{t+1}\le(1-2m\eta_t+(1+\rho)L^2\eta_t^2)v_t+\sigma^2\eta_t^2.
\]
The preceding scalar comparison, with $2mc>1$, gives
\begin{equation}\label{eq:localized_remainder}
v_t=O_N(t^{-1}),\qquad
\mathbb E[\|r(e_t)\|\mathbf1_{\{\tau_N>t\}}]=O_N(t^{-1}).
\end{equation}
Here constants may depend on the fixed starting index $N$.
We first take the iteration horizon to infinity and then let $N\to\infty$
in \eqref{eq:localization_probability}. Thus a remainder that vanishes
on these localized paths is $o_p(1)$ for the original process.

\paragraph{Matrix products.}
For $b\ge a$, let $P_{a,b}=\prod_{j=a}^b(I-\eta_jH)$, and set
$P_{a,b}=I$ when $a>b$. Choose $N$ large enough that
$\eta_j\|H\|<1$ for $j\ge N$. Spectral decomposition gives
\begin{equation}\label{eq:product_bounds}
\begin{aligned}
\|P_{s+1,t-1}\|&\le
\exp\left(-h\sum_{j=s+1}^{t-1}\eta_j\right), &&\alpha<1,\\
\|P_{s+1,t-1}\|&\le C_\nu
\left(\frac{s+\gamma}{t+\gamma}\right)^\nu,
&&\alpha=1,\quad 0<\nu<ch,
\end{aligned}
\end{equation}
for $t>s\ge N$. For the second bound, apply
$\log(1-c\lambda/(j+\gamma))=-c\lambda/(j+\gamma)+O(j^{-2})$
to each eigenvalue $\lambda$ of $H$ and sum over $j$.

\paragraph{A common martingale CLT argument.}
Equation~\eqref{eq:higher_noise_moments}, almost-sure convergence, and covariance continuity imply, for some
$\delta>0$,
\[
\mathbb E_t\xi_t=0,\qquad
\sup_t\mathbb E_t\|\xi_t\|^{2+\delta}\le C_\delta(\omega)<\infty,\qquad
S(\theta_t)\longrightarrow S_*\quad\text{almost surely}.
\]
Let $D_{T,s}$ be deterministic matrix weights such that
\begin{equation}\label{eq:weight_conditions}
\max_s\|D_{T,s}\|\to0,\qquad
\sup_T\sum_s\|D_{T,s}\|^2<\infty,\qquad
\sum_sD_{T,s}S_*D_{T,s}^\top\to\Sigma.
\end{equation}
Splitting the sum at a fixed index and then sending that index to infinity
shows that replacing $S_*$ by $S(\theta_s)$ changes the last limit by
$o(1)$ almost surely. For each $\epsilon>0$, conditional Lindeberg follows
from
\[
\sum_s\mathbb E_s\!\left[
\|D_{T,s}\xi_s\|^2\mathbf1_{\{\|D_{T,s}\xi_s\|>\epsilon\}}\right]
\le C_\delta(\omega)\epsilon^{-\delta}
\max_s\|D_{T,s}\|^\delta\sum_s\|D_{T,s}\|^2\longrightarrow0\quad\text{almost surely}.
\]
The martingale array CLT and the Cram\'er--Wold argument therefore give
$\sum_sD_{T,s}\xi_s\Rightarrow N(0,\Sigma)$
\cite[Chapter~3]{hallheyde1980}. This verifies explicitly the probabilistic
conditions used in the following two proofs.

\subsection{Last-iterate fluctuations}
\begin{proof}[Proof of Theorem~\ref{thm:last_clt}]
The exact recursion, started at a fixed sufficiently large $N$, is
\begin{equation}\label{eq:linear_representation}
e_t=P_{N,t-1}e_N-
\sum_{s=N}^{t-1}\eta_sP_{s+1,t-1}\xi_s-
\sum_{s=N}^{t-1}\eta_sP_{s+1,t-1}r(e_s).
\end{equation}
We verify directly the linearization and covariance limits underlying
the stochastic approximation CLT of Fabian \cite{fabian1968}.

If $\alpha<1$, the first term multiplied by $t^{\alpha/2}$ vanishes
by \eqref{eq:product_bounds}. On localized paths, the expected norm of
the scaled last term is bounded by
\[
C t^{\alpha/2}\sum_{s=N}^{t-1}\eta_s
\exp\left(-h\sum_{j=s+1}^{t-1}\eta_j\right)s^{-\alpha}
=O(t^{-\alpha/2}).
\]
To see this, split at $t/2$. The earlier terms have an exponentially
small factor $e^{-c_1t^{1-\alpha}}$; for the later terms,
$s^{-\alpha}\le C t^{-\alpha}$ and the remaining weighted geometric
sum is bounded. Localization then makes this term $o_p(1)$.

Set $D_{t,s}=t^{\alpha/2}\eta_sP_{s+1,t-1}$.
The same product bound gives
$\max_s\|D_{t,s}\|=O(t^{-\alpha/2})$ and
$\sum_s\|D_{t,s}\|^2=O(1)$.
On the window $t-s=O(\eta_t^{-1})$, uniformly on bounded rescaled
intervals, $\eta_s/\eta_t\to1$ and
$P_{s+1,t-1}=e^{-H\eta_t(t-s)}+o(1)$.
Riemann summation, using $t^\alpha\eta_t\to c$, gives
\[
\sum_{s=N}^{t-1}D_{t,s}S_*D_{t,s}^\top
\longrightarrow c\int_0^\infty e^{-Hv}S_*e^{-Hv}\,dv.
\]
The omitted tail beyond $t-s>R/\eta_t$ is bounded by $Ce^{-c_2R}$
after the negligible range $s<t/2$ is removed. This justifies sending
$R\to\infty$. The common martingale argument proves the first CLT.

If $\alpha=1$, fix $\nu\in(1/2,\min\{ch,1\})$.
The scaled initial term is $O_p(t^{1/2-\nu})=o_p(1)$.
By \eqref{eq:localized_remainder}, the localized expected norm of the
scaled nonlinear term is at most
\[
C_N t^{1/2-\nu}\sum_{s=N}^{t-1}s^{\nu-2}
=O_N(t^{1/2-\nu})=o(1).
\]
Set $D_{t,s}=\sqrt t\,\eta_sP_{s+1,t-1}$. Then
$\|D_{t,s}\|\le C t^{-1/2}(s/t)^{\nu-1}$, so the first two
conditions in \eqref{eq:weight_conditions} hold. Uniformly for
$s/t$ in compact subsets of $(0,1]$,
\[
\sqrt t\,D_{t,s}\longrightarrow c\,u^{cH-I}
\quad\text{when }s/t\to u.
\]
The squared-weight sum over $N\le s\le\varepsilon t$ is bounded by
$C\varepsilon^{2\nu-1}$. Hence Riemann summation gives
\[
\sum_{s=N}^{t-1}D_{t,s}S_*D_{t,s}^\top
\longrightarrow c^2\int_0^1u^{cH-I}S_*u^{cH-I}\,du.
\]
The martingale argument proves the endpoint CLT.

Finally, the first integral covariance solves
$H\Sigma_c+\Sigma_cH=cS_*$. In the second integral, the substitution
$u=e^{-v}$ gives
$c^2\int_0^\infty e^{-(cH-I/2)v}S_*e^{-(cH-I/2)v}\,dv$,
which solves the stated endpoint equation. Positive definiteness of
$H$ and $cH-I/2$, respectively, gives uniqueness and positive
semidefiniteness of these solutions.
\end{proof}

\subsection{Polyak--Ruppert averaging}
\begin{proof}[Proof of Theorem~\ref{thm:pr_clt}]
Suppose first that $\alpha<1$. Summing
$He_t=(e_t-e_{t+1})/\eta_t-\xi_t-r(e_t)$ from $N$ to $T$ gives
\[
H\sum_{t=N}^Te_t=-\sum_{t=N}^T\xi_t-
\sum_{t=N}^Tr(e_t)+R_{N,T},
\]
where summation by parts yields
\[
R_{N,T}=\frac{e_N}{\eta_N}-\frac{e_{T+1}}{\eta_T}
+\sum_{t=N+1}^T(\eta_t^{-1}-\eta_{t-1}^{-1})e_t.
\]
Since $\mathbb E\|e_t\|=O(t^{-\alpha/2})$ and
$\eta_t^{-1}-\eta_{t-1}^{-1}=O(t^{\alpha-1})$,
\[
T^{-1/2}\mathbb E\|R_{N,T}\|
=O_N(T^{-1/2})+O(T^{(\alpha-1)/2})=o(1).
\]
The localized nonlinear sum, divided by $\sqrt T$, has expected norm
at most
\[
C T^{-1/2}\sum_{t=N}^Tt^{-\alpha}=O(T^{1/2-\alpha})=o(1).
\]
Removing localization and the finite initial segment gives the classical
averaging representation
\cite{polyakjuditsky1992}
\[
\sqrt T\,\bar e_T=-H^{-1}\frac1{\sqrt T}\sum_{t=1}^T\xi_t+o_p(1).
\]
Taking $D_{T,s}=H^{-1}/\sqrt T$ in
\eqref{eq:weight_conditions} proves the first assertion.

For $\alpha=1$, use the product representation
\eqref{eq:linear_representation} and define
\[
B_{T,s}=\eta_s\sum_{t=s+1}^TP_{s+1,t-1},\qquad N\le s\le T,
\]
Then
\[
\sqrt T\,\bar e_T=-\frac1{\sqrt T}\sum_{s=N}^{T-1}B_{T,s}\xi_s
-\frac1{\sqrt T}\sum_{s=N}^{T-1}B_{T,s}r(e_s)+I_{N,T},
\]
where $B_{T,T}=0$ and $I_{N,T}$ contains the finite initial segment and the propagated
initial value. For $\nu\in(1/2,\min\{ch,1\})$,
\eqref{eq:product_bounds} implies
\begin{equation}\label{eq:averaging_weight_bound}
\|B_{T,s}\|\le C s^{\nu-1}\sum_{t=s+1}^Tt^{-\nu}
\le C(s/T)^{\nu-1},\qquad
I_{N,T}=O_p(T^{1/2-\nu}).
\end{equation}
By \eqref{eq:localized_remainder}, the localized nonlinear sum is
bounded in expected norm by
\[
\frac{C_N}{\sqrt T}\sum_{s=N}^{T-1}(s/T)^{\nu-1}s^{-1}
=O_N(T^{1/2-\nu})=o(1).
\]
Thus both remainders vanish in probability.

To identify the weights, the eigenvalue expansion used in
\eqref{eq:product_bounds} gives, uniformly when
$\varepsilon T\le s<t\le T$,
\[
P_{s+1,t-1}=((s+\gamma)/(t+\gamma))^{cH}+O_\varepsilon(T^{-1}).
\]
Consequently, uniformly on compact subsets of $(0,1]$,
\[
B_{T,\lfloor Tu\rfloor}\longrightarrow
K_c(u)=c\,u^{cH-I}\int_u^1v^{-cH}\,dv.
\]
Both $B_{T,s}$ and $K_c(s/T)$ satisfy the bound
$C(s/T)^{\nu-1}$. In particular,
\[
\frac1T\sum_{N\le s\le\varepsilon T}
\bigl(\|B_{T,s}\|^2+\|K_c(s/T)\|^2\bigr)
\le C\varepsilon^{2\nu-1}.
\]
Compact convergence and this bound establish
\begin{equation}\label{eq:kernel_l2}
\frac1T\sum_{s=N}^{T-1}\|B_{T,s}-K_c(s/T)\|^2\longrightarrow0.
\end{equation}
The relative variance bound, smoothness, and the strong-convexity MSE
estimate imply $\sup_s\mathbb E\|\xi_s\|^2<\infty$.
Martingale orthogonality therefore shows that
replacing $B_{T,s}$ by $K_c(s/T)$ in the noise sum produces an
$o_{L^2}(1)$ error. The finitely many terms with $s<N$ vanish because
$T^{-1/2}\|K_c(s/T)\|=O(T^{1/2-\nu})$ for fixed $s$.
We have therefore proved
\[
\sqrt T\,\bar e_T=-\frac1{\sqrt T}\sum_{s=1}^{T-1}K_c(s/T)\xi_s+o_p(1).
\]
For $D_{T,s}=T^{-1/2}K_c(s/T)$,
\eqref{eq:averaging_weight_bound} gives vanishing maximum weight and a
bounded squared-weight sum. The same compact/tail decomposition gives
\[
\sum_{s=1}^{T-1}D_{T,s}S_*D_{T,s}^\top
\longrightarrow\int_0^1K_c(u)S_*K_c(u)^\top\,du.
\]
The common martingale CLT argument completes the proof.
\end{proof}

\subsection{Nonconvex stationarity}
\begin{proof}
Smoothness, unbiasedness, and \eqref{eq:score_second_moment} give
\[
\mathbb E_tF(\theta_{t+1})
\le F(\theta_t)-\eta_t\left(1-\tfrac L2(1+\rho)\eta_t\right)
\|\nabla F(\theta_t)\|^2+\tfrac L2\sigma^2\eta_t^2.
\]
The stated step-size bound makes the descent coefficient at least
$\eta_t/2$. Taking expectations, summing, and using
$F(\theta_{T+1})\ge F_{\inf}$ proves the bound and the asserted rates.
\end{proof}


\subsection{Moment bounds for nonlinear-feature regression}
\label{app:regression}
For the regression model in Section~\ref{sec:experiments}, write
$Q=Z+U+iV$, $\theta^*=\theta_{\rm true}$, and
\[
A=\operatorname{Re}\{\psi(Q)\psi(Q)^\top\},\qquad
b=Y\operatorname{Re}\psi(Q),\qquad
H=\mathbb E[\psi(Z)\psi(Z)^\top].
\]
The corrected gradient is $G(\theta)=A\theta-b$.
Gaussian tails, the polynomial growth of $Y$, and the exponential growth
of the sine and cosine factors imply that $A$ and $b$ have moments of
every finite order. Conditional exact correction gives
$\mathbb EA=H$ and $\mathbb Eb=H\theta^*$.
The functions $\sin z$, $\cos(0.7z)$, and $z^2$ are linearly independent
on $\mathbb R$, and the latent density is positive everywhere.
Consequently $H\succ0$ and
\[
F(\theta)=\tfrac12(\theta-\theta^*)^\top H(\theta-\theta^*).
\]
Putting $e=\theta-\theta^*$ and $\mu=\lambda_{\min}(H)$, we obtain
\[
\xi(\theta)=(A-H)e+(A\theta^*-b),\qquad
\mathbb E\|\xi(\theta)\|^2
\le 2\mathbb E\|A\theta^*-b\|^2
+2\mathbb E\|A-H\|_{\rm op}^2\|e\|^2.
\]
Since $\|e\|\le\mu^{-1}\|\nabla F(\theta)\|$, this derives
\eqref{eq:relative_variance} directly from the data and score moments,
with finite constants. This argument verifies the required variance
condition for the joint response model, whose response-dependent envelope
is discussed in Section~\ref{sec:algorithm}. The proof of
Theorem~\ref{thm:strong_convex} therefore applies.
Here $\mu\approx0.382$ and $c=2.6$, so $2\mu c>1$ and the theoretical
MSE rate is $O(t^{-1})$.
The numerical value of $\mu$ is obtained from population quadrature.
